\chapter{Zadatak 2 --- skup nesigurnosti i robusna formulacija}
\label{ch:z2-nesigurnost}

\section{Što se ovdje radi}

Prevodi se rečenica ,,koeficijenti su nesigurni unutar $\pm15\%$`` u matematički
zapis, i formulira se problem koji traži dopustivost za \emph{sve} realizacije.

\begin{predznanje}
\begin{itemize}[nosep,leftmargin=*]
  \item Matrice $Q$, $c$, $G$, $h$ iz Zadatka 1 --- \eqref{eq:z1-Qc}, \eqref{eq:z1-Gh}.
  \item Poluprostor $a\trans x \leq b$ i njegovo crtanje --- \S\ref{sec:poluprostor}.
  \item Značenje $\sup$ nasuprot $\max$ --- \S\ref{sec:oznake-optimizacije}.
  \item Konveksan skup i konveksna ljuska --- \S\ref{sec:z1a-konveksan-skup}.
\end{itemize}
\end{predznanje}

\section{Intervali za pojedini koeficijent}
\label{sec:z2-intervali}

Nominalne vrijednosti iz Zadatka 1 su
\begin{equation}
a_{11}^{0} = -1, \quad a_{12}^{0} = -1, \quad a_{21}^{0} = 1, \quad a_{22}^{0} = -1 .
\end{equation}
Zadatak kaže da se stvarna vrijednost nalazi u intervalu od $0{,}85\,a_{ij}^{0}$
do $1{,}15\,a_{ij}^{0}$:
\begin{equation}
a_{ij} \in \bigl[\,0{,}85\,a_{ij}^{0},\ 1{,}15\,a_{ij}^{0}\,\bigr],
\qquad i,j \in \{1,2\}.
\label{eq:z2-interval}
\end{equation}

\begin{zamka}[Zamka: negativna nominalna vrijednost okreće poredak granica]
Za $a_{11}^{0} = -1$ doslovno uvrštavanje daje ,,od $-0{,}85$ do $-1{,}15$`` ---
a to nije interval u standardnom zapisu, jer je lijeva granica \emph{veća} od
desne. Ispravno je
\[
a_{11} \in [-1{,}15,\ -0{,}85].
\]
Množenje s negativnim brojem preokreće nejednakost, pa faktor $1{,}15$ daje
\emph{donju}, a $0{,}85$ \emph{gornju} granicu. Ako se to previdi, u nastavku se
za najgori slučaj bira kriva granica i cijelo rješenje ispada pogrešno.
\end{zamka}

\begin{formula}{Intervali koeficijenata}
\begin{equation}
a_{11},\, a_{12},\, a_{22} \in [-1{,}15,\ -0{,}85],
\qquad
a_{21} \in [0{,}85,\ 1{,}15] .
\label{eq:z2-intervali}
\end{equation}
Samo $a_{21}$ ima pozitivnu nominalnu vrijednost, pa je jedini čiji interval
ostaje ,,u očekivanom poretku``.
\end{formula}

\section{Skup nesigurnosti je pravokutnik}
\label{sec:z2-pravokutnik}

Retke matrice ograničenja označimo
\begin{equation}
a_1 = \begin{bmatrix} a_{11} \\ a_{12} \end{bmatrix},
\qquad
a_2 = \begin{bmatrix} a_{21} \\ a_{22} \end{bmatrix}.
\end{equation}
Kako svaki koeficijent varira nezavisno, svaki $a_i$ se kreće po
\emph{Kartezijevom produktu} svojih intervala:
\begin{formula}{Skupovi nesigurnosti}
\begin{align}
\mathcal{P}_1 &= [-1{,}15,\ -0{,}85] \times [-1{,}15,\ -0{,}85], \\
\mathcal{P}_2 &= [\phantom{-}0{,}85,\ \phantom{-}1{,}15] \times [-1{,}15,\ -0{,}85] .
\label{eq:z2-P}
\end{align}
Oba su pravokutnici u ravnini koeficijenata --- dakle poliedri, i to omeđeni, s
po četiri vrha.
\end{formula}

\begin{zamka}[Zamka: dva ograničenja imaju \emph{odvojene} skupove nesigurnosti]
Nesigurnost prvog retka nema veze s nesigurnošću drugog. Ne postoji zajednički
skup $\mathcal{P}$ za cijelu matricu; postoje dva neovisna, $\mathcal{P}_1$ i
$\mathcal{P}_2$. Zato se i robusni uvjet postavlja za svako ograničenje zasebno,
a ne za matricu odjednom. Da su koeficijenti bili povezani (npr.\ zajednički
faktor skaliranja), skup bi bio manji i rješenje manje konzervativno.
\end{zamka}

\section{Robusna formulacija}
\label{sec:z2-formulacija}

\begin{izvor}{Predavanje 8 --- Robusno optimiranje, slajd 2/21 (PDF str.~2)}
\slajd[0.86]{p08-robusan-lp}

\textbf{Transkripcija.} Promotrimo LP
\[
\min_x \ c\trans x \quad \text{uz ogr.} \quad a_i\trans x \leq b_i, \quad i = 1,\dots,m
\]
i pretpostavimo da su vektori $a_i \in \R^n$ nesigurni u smislu da ne znamo
njihove vrijednosti već znamo samo da $a_i \in \mathcal{P}_i$, gdje je
$\mathcal{P}_i = \{a \mid C_i a \leq d_i\}$ za poznate $C_i \in \R^{m_i \times n}$,
$d_i \in \R^{m_i}$. Pretpostavka je da skupovi (poliedri) $\mathcal{P}_i$ nisu
prazni.

\emph{Robusan LP:}
\[
\min_x \ c\trans x \quad \text{uz ogr.} \quad a_i\trans x \leq b_i \ \text{ za sve } a_i \in \mathcal{P}_i, \quad i = 1,\dots,m
\]

\textbf{Dvije stvari koje treba prenijeti na naš zadatak:}
\begin{itemize}[nosep,leftmargin=*]
  \item Nesigurnost je u \emph{koeficijentima ograničenja}, ne u funkciji cilja
        ni u desnim stranama --- točno kao kod nas.
  \item Riječi ,,za sve $a_i \in \mathcal{P}_i$`` pretvaraju jedno ograničenje u
        \emph{beskonačno mnogo} ograničenja, po jedno za svaku točku pravokutnika.
        Cijeli je posao svesti ih na konačno mnogo.
  \item Predavanje govori o LP-u; naša funkcija cilja je kvadratna. To ne smeta:
        nesigurnost je samo u ograničenjima, pa se svi zaključci o ograničenjima
        prenose nepromijenjeni, a funkcija cilja ostaje ista.
\end{itemize}
\end{izvor}

\begin{formula}{Robusni problem Zadatka 2}
\begin{equation}
\min_{x} \ \tfrac12 x\trans Q x + c\trans x
\quad \text{uz} \quad
a_i\trans x \leq h_i \ \text{ za sve } a_i \in \mathcal{P}_i, \ i = 1,2,
\label{eq:z2-robusni}
\end{equation}
gdje su $h_1 = -10$ i $h_2 = 3$ nepromijenjene desne strane.
\end{formula}

\section{Od ,,za sve`` do supremuma}
\label{sec:z2-supremum}

\begin{korak}[Prvi ključni korak]
Uvjet ,,$a_i\trans x \leq h_i$ za \emph{svaki} $a_i \in \mathcal{P}_i$``
ekvivalentan je jednom jedinom uvjetu na \emph{najgoru} vrijednost:
\begin{equation}
a_i\trans x \leq h_i \ \ \forall a_i \in \mathcal{P}_i
\qquad\Longleftrightarrow\qquad
\sup_{a_i \in \mathcal{P}_i} a_i\trans x \leq h_i .
\label{eq:z2-sup}
\end{equation}
Obrazloženje u oba smjera:
\begin{itemize}[nosep,leftmargin=*]
  \item ($\Rightarrow$) Ako svi članovi skupa zadovoljavaju nejednakost, onda je
        $h_i$ gornja granica tog skupa vrijednosti, pa je i \emph{najmanja}
        gornja granica --- supremum --- najviše $h_i$.
  \item ($\Leftarrow$) Ako je supremum $\leq h_i$, onda je i svaka pojedina
        vrijednost $\leq h_i$, jer nijedna ne premašuje supremum.
\end{itemize}
\end{korak}

Beskonačno mnogo nejednakosti time je zamijenjeno jednom --- ali cijena je što ta
jedna sadrži optimizacijski problem u sebi. Sljedeće poglavlje pokazuje da se za
poliedarski skup nesigurnosti taj unutarnji problem rješava egzaktno, i to samo
gledanjem vrhova.

\begin{izgovor}
,,Robusni problem tražim tako da ograničenje mora vrijediti za svaku realizaciju
koeficijenata iz zadanog pravokutnika. To je beskonačno mnogo nejednakosti, pa ih
sažimam u jednu: supremum od a i transponirano iks, po a i iz skupa nesigurnosti,
mora biti manji ili jednak desnoj strani. Ekvivalencija vrijedi u oba smjera ---
supremum je najmanja gornja granica, pa je uvjet na njega isto što i uvjet na
svaki član.``
\end{izgovor}
